\chapter{Geometric Numerical Integration}[Structure-Preserving Algorithms for Ordinary Differential Equations]

\section{Symplectic Integrators of Hamiltonian systems}
\subsection{Variational Integrators}
\subsubsection{Hamilton's Principle}
The motion extremizes the action integral is Hamilton's Principle

One can then derive the equation of the derivative of the action
\[
dS = p_1 dq_1 - p_0 dq_0
\]
which is the basic formula for simplectically generating functions.

\subsubsection{Discretization of Hamilton's principle}

I don't think I need to go over this

\subsubsection{Noether's theorem}
Continuous symmetris in te Lagrandian lead to first integrals, and thus this connection between Lagrandian and Hamiltonian leads to more first integrals.

\subsection{Characterization of Symplectic Methods}
Also trivial

\section{Non-Canonical Hamiltonian}

\subsection{Constrained Mechanical Systems}

For the most part these are pretty simple extensions of regular symplectic integrators

\subsection{Poisson System}

\subsubsection{Poisson Integrators}
